\bookchapter{Chapitre 10 : Paiement mobile et services du quotidien}

\chapterintro{Honnêtement, la génération de mes parents résistait au début à l'idée de « payer avec le smartphone ». « L'argent est dans le smartphone, que se passe-t-il si on se le fait voler ? » Cette inquiétude est normale, mais la réalité est : la probabilité de perdre les espèces dans votre poche est bien plus élevée que celle de se faire frauder en paiement mobile. Le paiement mobile ne stocke pas l'argent dans le smartphone, c'est juste une clé, une clé qui ouvre votre compte bancaire.}

\booksection{Comment lier une carte bancaire ?}

Que vous utilisiez Alipay ou WeChat, la première liaison de carte passe par ces étapes :

\begin{enumerate}
  \item Ouvrez l'application, trouvez « Moi » ou « Portefeuille ». Cliquez sur « Ajouter une carte bancaire ».
  \item Saisissez le numéro de carte. La plupart des applications permettent maintenant de scanner directement le numéro de carte, sans taper les chiffres un par un.
  \item La banque envoie un SMS avec un code de vérification au numéro enregistré. Saisissez-le, puis définissez un code de paiement à 6 chiffres.
  \item Ce code de paiement, le code PIN de votre carte bancaire et le code de verrouillage d'écran devraient idéalement être trois codes différents, ne les mettez pas identiques pour gagner du temps.
\end{enumerate}

Une petite astuce : si vous êtes préoccupé par la sécurité, vous pouvez ouvrir un compte électronique de catégorie II ou III dans l'application bancaire, dédié à la liaison avec le paiement mobile. Ne mettez sur cette carte que l'argent des petites dépenses courantes, même en cas de problème, les pertes sont limitées.

\booksection{Le code de paiement : « celui qui scanne encaisse »}

Beaucoup de gens confondent le code de paiement et le code de réception, cette distinction est importante.

\begin{itemize}
  \item \textbf{Code de paiement}: Le code-barres + QR code que vous montrez au caissier. Ce code équivaut à votre « portefeuille électronique », celui qui le scanne peut débiter votre compte. Donc\textbf{ne faites jamais de capture d'écran à envoyer à qui que ce soit}. Ce code se rafraîchit automatiquement chaque minute, justement pour éviter le vol.
  \item \textbf{Code de réception}: Les autres scannent ce code pour vous transférer de l'argent. Ce code peut être capturé et sauvegardé, envoyé à d'autres, sans risque.
\end{itemize}

Retenez bien cette différence, elle vous évitera bien des pièges d'arnaque. Toute personne qui vous demande « envoyez-moi une capture d'écran de votre code de paiement » est à 100 \% un escroc.

\booksection{Ces paramètres de sécurité, 5 minutes pour les configurer vous éviteront bien des soucis}

\begin{itemize}
  \item \textbf{Paiement par empreinte/visage}: Dans les paramètres de paiement d'Alipay et WeChat, activez le paiement biométrique. Ainsi, chaque paiement nécessite une empreinte ou un scan facial, même si vous perdez votre smartphone, celui qui le trouve ne peut pas dépenser votre argent.
  \item \textbf{Verrou de sécurité}: Alipay et WeChat ont tous deux une fonction « Verrou de sécurité ». Après activation, l'accès à la page « Portefeuille » de l'application nécessite une vérification d'identité supplémentaire. Chemin : Alipay Moi Paramètres Compte et sécurité Centre de sécurité Verrou de sécurité.
  \item \textbf{Vérification des paiements sans code}: Beaucoup de gens ont souscrit divers paiements sans code et l'ont oublié. Allez dans les « Paramètres de paiement » d'Alipay et WeChat, trouvez « Paiement sans code / Débit automatique », voyez quels services sont souscrits et désactivez ceux dont vous n'avez plus besoin.
  \item \textbf{Gestion des appareils connectés}: Vérifiez régulièrement quels appareils se sont connectés à votre compte. Si vous trouvez un appareil inconnu, déconnectez-le immédiatement. Chemin : Alipay Paramètres Compte et sécurité Gestion des appareils connectés.
\end{itemize}

\booksection{Bus et métro : le smartphone est votre carte de transport}

Maintenant, pour prendre le bus ou le métro, peu de gens utilisent encore une carte de transport physique. Pour passer les portiques avec le smartphone, il y a deux options :

\begin{itemize}
  \item \textbf{Carte de transport virtuelle dans le portefeuille du smartphone}: Ouvrez une carte de transport virtuelle directement dans Apple Wallet, Huawei Wallet, Xiaomi Wallet. L'avantage de cette solution : pas besoin de déverrouiller le smartphone pour passer, il suffit de l'approcher, la vitesse de réponse est la plus rapide, et ça fonctionne même sans réseau. Certaines villes offrent les mêmes avantages de correspondance que la carte physique.
  \item \textbf{Code de transport Alipay/WeChat}: Recherchez « Code de transport » dans l'application, activez le code de bus ou métro de votre ville. L'avantage : pas besoin de recharger, le montant est directement débité du solde ou de la carte bancaire liée. Mais nécessite une connexion réseau, et il faut ouvrir l'application pour trouver le code avant de passer.
\end{itemize}

Je recommande personnellement la carte de transport virtuelle du portefeuille du smartphone. Bien qu'il faille recharger 10 ou 20 yuans la première fois, c'est tellement plus pratique à l'usage, surtout quand on se précipite pour le métro, pas besoin de s'arrêter pour déverrouiller et chercher l'application.

\booksection{Paiement de factures, livraison, VTC}

\begin{itemize}
  \item \textbf{Paiement eau/électricité/gaz}: Dans « Paiement de factures » d'Alipay ou WeChat, liez les numéros de compte d'eau, d'électricité et de gaz de votre domicile. Ensuite, à chaque réception de facture, cliquez directement pour payer. Le numéro de compte se trouve généralement sur la facture papier de votre domicile, ou appelez le service client.
  \item \textbf{Comparaison de prix de livraison}: Le même restaurant a souvent des prix et des promotions différents sur Meituan et Ele.me. Regardez des deux côtés avant de commander, parfois vous économisez plusieurs yuans. De plus, la fonction « Pin Hao Fan » de Meituan, en commandant avec des amis, le prix peut être nettement réduit.
  \item \textbf{Sécurité VTC}: Dans Didi ou Gaode Dida, n'oubliez pas de configurer « Contact d'urgence » et « Partage de trajet ». Pendant la course, l'application partage automatiquement votre trajet avec le contact configuré. De plus, pour les courses de nuit, privilégiez « Taxi » plutôt que « Voiture rapide », les taxis sont soumis à une réglementation professionnelle, avec un coefficient de sécurité plus élevé.
\end{itemize}

\booksection{Gérer l'assurance maladie, la sécurité sociale et les démarches administratives avec le smartphone}

Les changements ces dernières années sont considérables, beaucoup de démarches qui nécessitaient auparavant de faire la queue dans un guichet se font maintenant en quelques minutes sur le smartphone.

\begin{itemize}
  \item \textbf{Certificat électronique d'assurance maladie}: Recherchez « Certificat électronique d'assurance maladie » dans Alipay ou WeChat, après activation, passez simplement votre smartphone à l'hôpital ou à la pharmacie, sans carte physique. La première activation nécessite une vérification par reconnaissance faciale, ensuite un code dynamique est généré à chaque utilisation.
  \item \textbf{Consultation de la sécurité sociale}: Téléchargez l'application « Zhangshang 12333 » ou utilisez le « Centre citoyen » dans Alipay, vous pouvez consulter les cotisations de sécurité sociale, le solde du compte de retraite, les demandes d'allocation chômage, etc.
  \item \textbf{Impôt sur le revenu}: La déclaration annuelle de mars à juin se fait sur l'application « Impôt sur le revenu personnel ». Vérifiez si vous pouvez déclarer des déductions supplémentaires spéciales (éducation des enfants, intérêts de prêt immobilier, prise en charge des personnes âgées, etc.), qui peuvent vous aider à payer moins d'impôts ou obtenir un remboursement.
\end{itemize}

\begin{tcbwarn}{Retenez ces règles, pour vous protéger des arnaques}
\begin{itemize}
  \item Les autorités judiciaires, les banques et le service client ne vous demanderont jamais de transférer vers un « compte sécurisé ».
  \item Toute personne vous demandant d'envoyer une capture d'écran de votre code de paiement est un escroc.
  \item Ne cliquez pas sur les liens inconnus, surtout les SMS du type « ETC expiré » ou « Carte de sécurité sociale suspendue ».
  \item Les remboursements d'achats en ligne reviennent par le même chemin, pas besoin de fournir votre numéro de carte bancaire.
\end{itemize}

\end{tcbwarn}

\begin{tcbpractice}{Faites-le maintenant : effectuez un bilan de sécurité de vos applications de paiement}
\begin{enumerate}
  \item Ouvrez Alipay et WeChat, confirmez que le paiement par empreinte/visage est activé.
  \item Trouvez et activez la fonction « Verrou de sécurité ».
  \item Vérifiez la liste « Paiement sans code / Débit automatique », désactivez les services dont vous n'avez plus besoin.
  \item Consultez la liste « Appareils connectés », retirez les appareils inconnus.
  \item Activez votre certificat électronique d'assurance maladie, plus besoin de carte physique à votre prochain passage à l'hôpital.
\end{enumerate}

\end{tcbpractice}
